\documentclass[11pt,a4paper]{article}

\newcommand{\doctitle}{Mortgage Pricing Model (MKM-MP-001)}
\newcommand{\docsubtitle}{Wolfram Oracle Benchmark of the Fair-Value Calculation, after the Fix}
\newcommand{\docversion}{0.2}
\newcommand{\docdate}{27-October-2026}
\newcommand{\docauthor}{David K Kelly}

% MKM Research Labs — Shared Document Header
% Usage: % MKM Research Labs — Shared Document Header
% Usage: % MKM Research Labs — Shared Document Header
% Usage: \input{../shared/mkm_header}  (before \begin{document})
%
% Required commands before \input:
%   \newcommand{\doctitle}{...}
%   \newcommand{\docsubtitle}{...}
%   \newcommand{\docversion}{1.0}
%   \newcommand{\docdate}{13-February-2026}
%   \newcommand{\docauthor}{David K Kelly}

% ─── Packages ───────────────────────────────────────────────────────────────
\usepackage[utf8]{inputenc}
\usepackage{amsmath,amsfonts,amssymb}
\usepackage{listings}
\usepackage{color,xcolor}
\usepackage{graphicx}
\usepackage{hyperref}
\usepackage[toc,page]{appendix}
\usepackage{tabto}
\usepackage{float}
\usepackage{comment}
\usepackage[en-GB]{datetime2}
\usepackage{setspace}
\usepackage{subcaption}
\usepackage{booktabs}
\usepackage{longtable}
\usepackage{adjustbox}
\usepackage{enumitem}
\usepackage{fancyhdr}
\usepackage{lastpage}

% ─── Page geometry ──────────────────────────────────────────────────────────
\onehalfspacing
\usepackage{geometry}
\geometry{a4paper, total={170mm,247mm}, left=20mm, top=25mm, bottom=25mm}

% ─── Headers and footers ───────────────────────────────────────────────────
\pagestyle{fancy}
\fancyhf{}
\fancyhead[L]{\small\textcolor{gray}{MKM Research Labs}}
\fancyhead[R]{\small\textcolor{gray}{\doctitle{} --- v\docversion}}
\fancyfoot[C]{\small\textcolor{gray}{CONFIDENTIAL}}
\fancyfoot[R]{\small\textcolor{gray}{Page \thepage\ of \pageref{LastPage}}}
\renewcommand{\headrulewidth}{0.4pt}
\renewcommand{\footrulewidth}{0.4pt}

% ─── Colours ────────────────────────────────────────────────────────────────
\definecolor{mkmblue}{HTML}{1565C0}
\definecolor{codegreen}{rgb}{0,0.6,0}
\definecolor{codegray}{rgb}{0.5,0.5,0.5}
\definecolor{codepurple}{rgb}{0.58,0,0.82}
\definecolor{backcolour}{rgb}{0.95,0.95,0.92}

% ─── Hyperref setup ─────────────────────────────────────────────────────────
\hypersetup{
    colorlinks=true,
    linkcolor=mkmblue,
    citecolor=mkmblue,
    urlcolor=mkmblue
}

% ─── Code listings ──────────────────────────────────────────────────────────
\lstdefinestyle{mystyle}{
    backgroundcolor=\color{backcolour},
    commentstyle=\color{codegreen},
    keywordstyle=\color{magenta},
    numberstyle=\tiny\color{codegray},
    stringstyle=\color{codepurple},
    basicstyle=\ttfamily\footnotesize,
    breakatwhitespace=false,
    breaklines=true,
    captionpos=b,
    keepspaces=true,
    numbers=left,
    numbersep=5pt,
    showspaces=false,
    showstringspaces=false,
    showtabs=false,
    tabsize=2
}
\lstset{style=mystyle}

% ─── Section formatting ────────────────────────────────────────────────────
\setcounter{tocdepth}{3}
\setcounter{secnumdepth}{3}

% ─── Shared macros ──────────────────────────────────────────────────────────
\newcommand{\mkmtitlepage}{%
    \begin{titlepage}
    \begin{center}
    \vspace*{1cm}
    \begin{figure}[H]\centering
        \includegraphics[width=0.4\textwidth]{../shared/Logos/MKM_logo.png}
    \end{figure}
    \vspace{1.5cm}
    {\Huge\bfseries\textcolor{mkmblue}{\doctitle}\par}
    \vspace{0.5cm}
    {\Large\textcolor{gray}{\docsubtitle}\par}
    \vspace{2cm}
    {\large MKM Research Labs\par}
    \vspace{0.5cm}
    {\large Version \docversion{} --- \docdate\par}
    \vspace{0.5cm}
    {\large \docauthor\par}
    \vspace{2cm}
    \textcolor{red}{\textbf{CONFIDENTIAL --- PROPRIETARY}}
    \end{center}
    \end{titlepage}
}

\newcommand{\mkmlegalpage}{%
    \clearpage
    \thispagestyle{empty}
    \begin{center}
    \large\textbf{Legal Notice}
    \vspace{1em}

    \textbf{PROPRIETARY AND CONFIDENTIAL}
    \end{center}

    \noindent This document, including all algorithms, models, methodology, formulas,
    calculations, and intellectual concepts contained herein, constitutes the exclusive
    intellectual property and confidential trade secrets of MKM Research Labs (``MKM'').

    \vspace{0.5em}
    \noindent Any usage, reproduction, distribution, modification, or disclosure of this
    document or any portion thereof, without the express prior written authorisation from
    MKM Research Labs is strictly prohibited and constitutes an infringement of intellectual
    property rights.

    \vspace{0.5em}
    \noindent The document is provided on an ``as is'' basis. MKM Research Labs makes no
    representations or warranties of any kind, express or implied, regarding its accuracy,
    reliability, or suitability for any particular purpose.

    \vspace{1em}
    \begin{center}
    All rights reserved. \copyright\ 2021--2026 MKM Research Labs.

    \vspace{0.5em}
    Governed by the laws of the United Kingdom.
    \end{center}
    \clearpage
}

\newcommand{\mkmhistorypage}[1]{%
    \section{Document History}
    \begin{table}[ht]
        \centering
        \begin{tabular}{l|l|l|l}
            \textbf{Date} & \textbf{Version} & \textbf{Description} & \textbf{Author} \\
            \hline
            #1
        \end{tabular}
        \label{tab:revision_history}
    \end{table}
    \clearpage
}

% ─── Bank validation protocol section commands ──────────────────────────────
% Use these to ensure every model doc has the required sections.
% They are just styled \section calls — content is provided by each doc.

\newcommand{\mkmexecsummary}{\section{Executive Summary}}
\newcommand{\mkmpurpose}{\section{Model Purpose and Scope}}
\newcommand{\mkmframework}{\section{Mathematical Framework}}
\newcommand{\mkminputs}{\section{Input Parameters and Calibration}}
\newcommand{\mkmimplementation}{\section{Implementation Details}}
\newcommand{\mkmsensitivity}{\section{Sensitivity Analysis}}
\newcommand{\mkmvalidation}{\section{Validation and Backtesting}}
\newcommand{\mkmtestresults}{\subsection{Automated Test Results}\input{test_results}}
\newcommand{\mkmlimitations}{\section{Model Limitations and Known Weaknesses}}

\newcommand{\mkmgovernance}{\section{Governance Validation Questions}
\noindent The following nine questions are mandated by Chapter~5 of the
\emph{Model Risk Governance for Vendors} handbook and must be explicitly
addressed for every model in the inventory. Response status is tracked in the
Model Governance Dashboard (MRC portal). \par\medskip}

\newcommand{\mkmvqitem}[3]{\subsection{Q#1 --- #2}
\noindent\textit{``#3''}\par\medskip\noindent}

\newcommand{\mkmchangelog}{\section{Change History}}
  (before \begin{document})
%
% Required commands before \input:
%   \newcommand{\doctitle}{...}
%   \newcommand{\docsubtitle}{...}
%   \newcommand{\docversion}{1.0}
%   \newcommand{\docdate}{13-February-2026}
%   \newcommand{\docauthor}{David K Kelly}

% ─── Packages ───────────────────────────────────────────────────────────────
\usepackage[utf8]{inputenc}
\usepackage{amsmath,amsfonts,amssymb}
\usepackage{listings}
\usepackage{color,xcolor}
\usepackage{graphicx}
\usepackage{hyperref}
\usepackage[toc,page]{appendix}
\usepackage{tabto}
\usepackage{float}
\usepackage{comment}
\usepackage[en-GB]{datetime2}
\usepackage{setspace}
\usepackage{subcaption}
\usepackage{booktabs}
\usepackage{longtable}
\usepackage{adjustbox}
\usepackage{enumitem}
\usepackage{fancyhdr}
\usepackage{lastpage}

% ─── Page geometry ──────────────────────────────────────────────────────────
\onehalfspacing
\usepackage{geometry}
\geometry{a4paper, total={170mm,247mm}, left=20mm, top=25mm, bottom=25mm}

% ─── Headers and footers ───────────────────────────────────────────────────
\pagestyle{fancy}
\fancyhf{}
\fancyhead[L]{\small\textcolor{gray}{MKM Research Labs}}
\fancyhead[R]{\small\textcolor{gray}{\doctitle{} --- v\docversion}}
\fancyfoot[C]{\small\textcolor{gray}{CONFIDENTIAL}}
\fancyfoot[R]{\small\textcolor{gray}{Page \thepage\ of \pageref{LastPage}}}
\renewcommand{\headrulewidth}{0.4pt}
\renewcommand{\footrulewidth}{0.4pt}

% ─── Colours ────────────────────────────────────────────────────────────────
\definecolor{mkmblue}{HTML}{1565C0}
\definecolor{codegreen}{rgb}{0,0.6,0}
\definecolor{codegray}{rgb}{0.5,0.5,0.5}
\definecolor{codepurple}{rgb}{0.58,0,0.82}
\definecolor{backcolour}{rgb}{0.95,0.95,0.92}

% ─── Hyperref setup ─────────────────────────────────────────────────────────
\hypersetup{
    colorlinks=true,
    linkcolor=mkmblue,
    citecolor=mkmblue,
    urlcolor=mkmblue
}

% ─── Code listings ──────────────────────────────────────────────────────────
\lstdefinestyle{mystyle}{
    backgroundcolor=\color{backcolour},
    commentstyle=\color{codegreen},
    keywordstyle=\color{magenta},
    numberstyle=\tiny\color{codegray},
    stringstyle=\color{codepurple},
    basicstyle=\ttfamily\footnotesize,
    breakatwhitespace=false,
    breaklines=true,
    captionpos=b,
    keepspaces=true,
    numbers=left,
    numbersep=5pt,
    showspaces=false,
    showstringspaces=false,
    showtabs=false,
    tabsize=2
}
\lstset{style=mystyle}

% ─── Section formatting ────────────────────────────────────────────────────
\setcounter{tocdepth}{3}
\setcounter{secnumdepth}{3}

% ─── Shared macros ──────────────────────────────────────────────────────────
\newcommand{\mkmtitlepage}{%
    \begin{titlepage}
    \begin{center}
    \vspace*{1cm}
    \begin{figure}[H]\centering
        \includegraphics[width=0.4\textwidth]{../shared/Logos/MKM_logo.png}
    \end{figure}
    \vspace{1.5cm}
    {\Huge\bfseries\textcolor{mkmblue}{\doctitle}\par}
    \vspace{0.5cm}
    {\Large\textcolor{gray}{\docsubtitle}\par}
    \vspace{2cm}
    {\large MKM Research Labs\par}
    \vspace{0.5cm}
    {\large Version \docversion{} --- \docdate\par}
    \vspace{0.5cm}
    {\large \docauthor\par}
    \vspace{2cm}
    \textcolor{red}{\textbf{CONFIDENTIAL --- PROPRIETARY}}
    \end{center}
    \end{titlepage}
}

\newcommand{\mkmlegalpage}{%
    \clearpage
    \thispagestyle{empty}
    \begin{center}
    \large\textbf{Legal Notice}
    \vspace{1em}

    \textbf{PROPRIETARY AND CONFIDENTIAL}
    \end{center}

    \noindent This document, including all algorithms, models, methodology, formulas,
    calculations, and intellectual concepts contained herein, constitutes the exclusive
    intellectual property and confidential trade secrets of MKM Research Labs (``MKM'').

    \vspace{0.5em}
    \noindent Any usage, reproduction, distribution, modification, or disclosure of this
    document or any portion thereof, without the express prior written authorisation from
    MKM Research Labs is strictly prohibited and constitutes an infringement of intellectual
    property rights.

    \vspace{0.5em}
    \noindent The document is provided on an ``as is'' basis. MKM Research Labs makes no
    representations or warranties of any kind, express or implied, regarding its accuracy,
    reliability, or suitability for any particular purpose.

    \vspace{1em}
    \begin{center}
    All rights reserved. \copyright\ 2021--2026 MKM Research Labs.

    \vspace{0.5em}
    Governed by the laws of the United Kingdom.
    \end{center}
    \clearpage
}

\newcommand{\mkmhistorypage}[1]{%
    \section{Document History}
    \begin{table}[ht]
        \centering
        \begin{tabular}{l|l|l|l}
            \textbf{Date} & \textbf{Version} & \textbf{Description} & \textbf{Author} \\
            \hline
            #1
        \end{tabular}
        \label{tab:revision_history}
    \end{table}
    \clearpage
}

% ─── Bank validation protocol section commands ──────────────────────────────
% Use these to ensure every model doc has the required sections.
% They are just styled \section calls — content is provided by each doc.

\newcommand{\mkmexecsummary}{\section{Executive Summary}}
\newcommand{\mkmpurpose}{\section{Model Purpose and Scope}}
\newcommand{\mkmframework}{\section{Mathematical Framework}}
\newcommand{\mkminputs}{\section{Input Parameters and Calibration}}
\newcommand{\mkmimplementation}{\section{Implementation Details}}
\newcommand{\mkmsensitivity}{\section{Sensitivity Analysis}}
\newcommand{\mkmvalidation}{\section{Validation and Backtesting}}
\newcommand{\mkmtestresults}{\subsection{Automated Test Results}\input{test_results}}
\newcommand{\mkmlimitations}{\section{Model Limitations and Known Weaknesses}}

\newcommand{\mkmgovernance}{\section{Governance Validation Questions}
\noindent The following nine questions are mandated by Chapter~5 of the
\emph{Model Risk Governance for Vendors} handbook and must be explicitly
addressed for every model in the inventory. Response status is tracked in the
Model Governance Dashboard (MRC portal). \par\medskip}

\newcommand{\mkmvqitem}[3]{\subsection{Q#1 --- #2}
\noindent\textit{``#3''}\par\medskip\noindent}

\newcommand{\mkmchangelog}{\section{Change History}}
  (before \begin{document})
%
% Required commands before \input:
%   \newcommand{\doctitle}{...}
%   \newcommand{\docsubtitle}{...}
%   \newcommand{\docversion}{1.0}
%   \newcommand{\docdate}{13-February-2026}
%   \newcommand{\docauthor}{David K Kelly}

% ─── Packages ───────────────────────────────────────────────────────────────
\usepackage[utf8]{inputenc}
\usepackage{amsmath,amsfonts,amssymb}
\usepackage{listings}
\usepackage{color,xcolor}
\usepackage{graphicx}
\usepackage{hyperref}
\usepackage[toc,page]{appendix}
\usepackage{tabto}
\usepackage{float}
\usepackage{comment}
\usepackage[en-GB]{datetime2}
\usepackage{setspace}
\usepackage{subcaption}
\usepackage{booktabs}
\usepackage{longtable}
\usepackage{adjustbox}
\usepackage{enumitem}
\usepackage{fancyhdr}
\usepackage{lastpage}

% ─── Page geometry ──────────────────────────────────────────────────────────
\onehalfspacing
\usepackage{geometry}
\geometry{a4paper, total={170mm,247mm}, left=20mm, top=25mm, bottom=25mm}

% ─── Headers and footers ───────────────────────────────────────────────────
\pagestyle{fancy}
\fancyhf{}
\fancyhead[L]{\small\textcolor{gray}{MKM Research Labs}}
\fancyhead[R]{\small\textcolor{gray}{\doctitle{} --- v\docversion}}
\fancyfoot[C]{\small\textcolor{gray}{CONFIDENTIAL}}
\fancyfoot[R]{\small\textcolor{gray}{Page \thepage\ of \pageref{LastPage}}}
\renewcommand{\headrulewidth}{0.4pt}
\renewcommand{\footrulewidth}{0.4pt}

% ─── Colours ────────────────────────────────────────────────────────────────
\definecolor{mkmblue}{HTML}{1565C0}
\definecolor{codegreen}{rgb}{0,0.6,0}
\definecolor{codegray}{rgb}{0.5,0.5,0.5}
\definecolor{codepurple}{rgb}{0.58,0,0.82}
\definecolor{backcolour}{rgb}{0.95,0.95,0.92}

% ─── Hyperref setup ─────────────────────────────────────────────────────────
\hypersetup{
    colorlinks=true,
    linkcolor=mkmblue,
    citecolor=mkmblue,
    urlcolor=mkmblue
}

% ─── Code listings ──────────────────────────────────────────────────────────
\lstdefinestyle{mystyle}{
    backgroundcolor=\color{backcolour},
    commentstyle=\color{codegreen},
    keywordstyle=\color{magenta},
    numberstyle=\tiny\color{codegray},
    stringstyle=\color{codepurple},
    basicstyle=\ttfamily\footnotesize,
    breakatwhitespace=false,
    breaklines=true,
    captionpos=b,
    keepspaces=true,
    numbers=left,
    numbersep=5pt,
    showspaces=false,
    showstringspaces=false,
    showtabs=false,
    tabsize=2
}
\lstset{style=mystyle}

% ─── Section formatting ────────────────────────────────────────────────────
\setcounter{tocdepth}{3}
\setcounter{secnumdepth}{3}

% ─── Shared macros ──────────────────────────────────────────────────────────
\newcommand{\mkmtitlepage}{%
    \begin{titlepage}
    \begin{center}
    \vspace*{1cm}
    \begin{figure}[H]\centering
        \includegraphics[width=0.4\textwidth]{../shared/Logos/MKM_logo.png}
    \end{figure}
    \vspace{1.5cm}
    {\Huge\bfseries\textcolor{mkmblue}{\doctitle}\par}
    \vspace{0.5cm}
    {\Large\textcolor{gray}{\docsubtitle}\par}
    \vspace{2cm}
    {\large MKM Research Labs\par}
    \vspace{0.5cm}
    {\large Version \docversion{} --- \docdate\par}
    \vspace{0.5cm}
    {\large \docauthor\par}
    \vspace{2cm}
    \textcolor{red}{\textbf{CONFIDENTIAL --- PROPRIETARY}}
    \end{center}
    \end{titlepage}
}

\newcommand{\mkmlegalpage}{%
    \clearpage
    \thispagestyle{empty}
    \begin{center}
    \large\textbf{Legal Notice}
    \vspace{1em}

    \textbf{PROPRIETARY AND CONFIDENTIAL}
    \end{center}

    \noindent This document, including all algorithms, models, methodology, formulas,
    calculations, and intellectual concepts contained herein, constitutes the exclusive
    intellectual property and confidential trade secrets of MKM Research Labs (``MKM'').

    \vspace{0.5em}
    \noindent Any usage, reproduction, distribution, modification, or disclosure of this
    document or any portion thereof, without the express prior written authorisation from
    MKM Research Labs is strictly prohibited and constitutes an infringement of intellectual
    property rights.

    \vspace{0.5em}
    \noindent The document is provided on an ``as is'' basis. MKM Research Labs makes no
    representations or warranties of any kind, express or implied, regarding its accuracy,
    reliability, or suitability for any particular purpose.

    \vspace{1em}
    \begin{center}
    All rights reserved. \copyright\ 2021--2026 MKM Research Labs.

    \vspace{0.5em}
    Governed by the laws of the United Kingdom.
    \end{center}
    \clearpage
}

\newcommand{\mkmhistorypage}[1]{%
    \section{Document History}
    \begin{table}[ht]
        \centering
        \begin{tabular}{l|l|l|l}
            \textbf{Date} & \textbf{Version} & \textbf{Description} & \textbf{Author} \\
            \hline
            #1
        \end{tabular}
        \label{tab:revision_history}
    \end{table}
    \clearpage
}

% ─── Bank validation protocol section commands ──────────────────────────────
% Use these to ensure every model doc has the required sections.
% They are just styled \section calls — content is provided by each doc.

\newcommand{\mkmexecsummary}{\section{Executive Summary}}
\newcommand{\mkmpurpose}{\section{Model Purpose and Scope}}
\newcommand{\mkmframework}{\section{Mathematical Framework}}
\newcommand{\mkminputs}{\section{Input Parameters and Calibration}}
\newcommand{\mkmimplementation}{\section{Implementation Details}}
\newcommand{\mkmsensitivity}{\section{Sensitivity Analysis}}
\newcommand{\mkmvalidation}{\section{Validation and Backtesting}}
\newcommand{\mkmtestresults}{\subsection{Automated Test Results}\input{test_results}}
\newcommand{\mkmlimitations}{\section{Model Limitations and Known Weaknesses}}

\newcommand{\mkmgovernance}{\section{Governance Validation Questions}
\noindent The following nine questions are mandated by Chapter~5 of the
\emph{Model Risk Governance for Vendors} handbook and must be explicitly
addressed for every model in the inventory. Response status is tracked in the
Model Governance Dashboard (MRC portal). \par\medskip}

\newcommand{\mkmvqitem}[3]{\subsection{Q#1 --- #2}
\noindent\textit{``#3''}\par\medskip\noindent}

\newcommand{\mkmchangelog}{\section{Change History}}

% T1 so the underscores in code identifiers print as underscores.
\usepackage[T1]{fontenc}
\usepackage{lmodern}
% Generated by tools.wolfram_oracle -- do not edit.
\newcommand{\bmGenerated}{2026-10-27}
\newcommand{\bmCodeCommit}{\texttt{2234867eb786}}
\newcommand{\bmCodeDirty}{}
\newcommand{\bmNumerics}{numpy 1.26.4, scipy 1.15.1, mpmath 1.4.1}
\newcommand{\bmWolfram}{15.0.0 for Mac OS X ARM (64-bit) (May 26, 2026)}
\newcommand{\bmWolframVia}{ via live kernel session (wolframclient)}
\newcommand{\bmWolframConfirms}{Wolfram and the 50-digit reference agree to $2.8\times10^{-30}$ on both values.}
\newcommand{\bmCases}{360}
\newcommand{\bmChecks}{6}
\newcommand{\bmOutcome}{all pass}
\newcommand{\bmCodeErr}{$1.7\times10^{-14}$}
\newcommand{\bmCorrection}{34.6}
\newcommand{\bmCorrectionShare}{53}
\newcommand{\bmWorstLoan}{300{,}000}
\newcommand{\bmWorstOld}{101{,}992}
\newcommand{\bmWorstNew}{205{,}812}
\newcommand{\bmPayLow}{8.3}
\newcommand{\bmPayHigh}{11.9}
\newcommand{\bmZeroIncome}{37.625}

% The shared history table sets every column to its natural width; wrap the description.
\renewcommand{\mkmhistorypage}[1]{%
    \section{Document History}
    \begin{table}[ht]
        \centering
        \begin{tabular}{l|l|p{9cm}|l}
            \textbf{Date} & \textbf{Version} & \textbf{Description} & \textbf{Author} \\
            \hline
            #1
        \end{tabular}
    \end{table}
}

\begin{document}
% Long code paths and identifiers cannot hyphenate; let the spacing stretch instead.
\sloppy

\mkmtitlepage
\mkmlegalpage

\tableofcontents
\clearpage

% =========================================================================
\mkmexecsummary
% =========================================================================

This document benchmarks the Mortgage Pricing Model (MKM-MP-001), which values a
flood-exposed residential mortgage from its amortisation schedule, an
affordability-based credit spread and a loss given default. Version~0.1 tested the
pricer against model document v1.1 and found that the fair value counted the expected
loss twice and took the recovery a month late. Both are fixed in the pricer, and the
document is corrected as v1.3, which this version tests.

\paragraph{The code computes its documented formulae.} Over \bmCases{} loans ---
loan-to-value 0.5 to 1.0, recovery haircuts 0 to 40\%, two incomes, two coupons,
three terms, discounted at the coupon and at a risk-free rate --- the payment, spread,
expected cashflows, expected losses and fair value agree with a 50-digit
implementation of v1.3's Eqs.~1--24 to \bmCodeErr. \bmWolframConfirms

\paragraph{The fair value counts each loss once.} The fair value is now the present
value of the expected cashflows, whose recovery is already net of the loss, and it
equals Wolfram's independent pathwise valuation: the expectation, over the month of
default, of the payments made before it and the recovery made in it. The loss is
taken on the balance owed when the month of default began, so the lender never
recovers more than the property fetches. The correction raises fair value on the
\bmCorrectionShare\% of loans in the grid that carry a loss given default, by up to
\bmCorrection\% of par: a GBP~\bmWorstLoan{} loan that v1.1 valued at
GBP~\bmWorstOld{} is worth GBP~\bmWorstNew{} (Table~\ref{tab:bm-correction}).
Expected loss itself is reported as before, now on the opening balance.

\paragraph{The document's claims hold} (Table~\ref{tab:bm-claims}). v1.3 corrects
those v0.1 found wrong: a 100\,bp rise raises a 25-year payment by
\bmPayLow--\bmPayHigh\%, within the 8--12\% now stated; missing income prices at the
top of the spread schedule (\bmZeroIncome\% at the baseline), as now stated, the
15\% fallback being unreachable; and discounting is described as production does it,
at the risk-free rate the loan calculator supplies. Table~7 is reproduced exactly,
the expected-loss bound holds, and Wolfram proves the par identity symbolically.

\paragraph{Evidence.} \bmChecks{} checks, outcome: \bmOutcome{}
(Table~\ref{tab:bm-checks}).
\begin{itemize}[leftmargin=*,nosep]
\item Wolfram: \bmWolfram{}\bmWolframVia.
\item Code under test, at commit \bmCodeCommit{}\bmCodeDirty.
\item Numerics: \bmNumerics.
\item Generated \bmGenerated{} by \texttt{python -m tools.wolfram\_oracle mp001}.
\end{itemize}

% =========================================================================
\mkmpurpose
% =========================================================================

\subsection{What is benchmarked}
\texttt{LoanPricer.price\_loan} in \texttt{src/models/loan/pricer/} (\texttt{\_pricing.py},
with the spread in \texttt{credit.py}), imported as production imports it and called
as the loan calculator (\texttt{src/routes/\_loan\_pricing}) calls it.

\subsection{What is not}
The calibration of the spread schedule (\S4.3), the coupon the calculator builds
upstream of the pricer, and the portfolio aggregation (\S3.9), which is arithmetic.
This is a benchmark of the implementation (validation question Q8), not independent
validation (Q9).

% =========================================================================
\mkmframework
% =========================================================================

Following the document (v1.3, \S3): monthly payment $M$, balances $B_0,\dots,B_n$,
monthly default probability $h=1-e^{-s/12}$, survival $S_i=(1-h)^i$,
$\mathrm{PD}_i=S_{i-1}-S_i$, recoverable value $R=(1-\delta)V$, discount factors $D_i$,
and the loss on a default in month $i$ taken on the balance owed when the month began,
$\mathrm{LGD}_i=\max(0,B_{i-1}-R)$. Then
\begin{equation}
\mathrm{FV}=\mathrm{PV}_{\mathrm{CF}}=\sum_i D_i\bigl[M S_i+(B_{i-1}-\mathrm{LGD}_i)\,\mathrm{PD}_i\bigr],
\qquad
\mathrm{PV}_{\mathrm{EL}}=\sum_i D_i\,\mathrm{LGD}_i\,\mathrm{PD}_i,
\end{equation}
with $B_{i-1}-\mathrm{LGD}_i=\min(B_{i-1},R)$, the recovery.

\paragraph{The pathwise value.} Let $\tau$ be the month of default, so that
$P(\tau=i)=\mathrm{PD}_i$ and $P(\tau>n)=S_n$. On the path where it defaults in month
$\tau$, the lender receives $M$ in each month before $\tau$ and $\min(B_{\tau-1},R)$
in month $\tau$; on the path where it never defaults, all $n$ payments. Taking the
expectation,
\begin{equation}
\mathbb{E}[\mathrm{PV}]=\sum_{i=1}^n D_i\bigl[M\,P(\tau>i)+\min(B_{i-1},R)\,P(\tau=i)\bigr]
=\mathrm{PV}_{\mathrm{CF}}.
\end{equation}

\paragraph{What v1.1 did.} v1.1 took the loss on the closing balance,
$\max(0,B_i-R)$, so its recovery $R+(B_{i-1}-B_i)$ exceeded $R$ by a month's principal
when the loan was under water; and it set $\mathrm{FV}=\mathrm{PV}_{\mathrm{CF}}-\mathrm{PV}_{\mathrm{EL}}$,
deducting a loss already inside the recovery.

% =========================================================================
\section{Benchmark Method}
% =========================================================================

The equations above are evaluated with \texttt{mpmath} at 50 significant digits, from
v1.3 of the document rather than the code; the loan-to-value ratio is taken as the double the
code compares against its bands, since a band edge such as 0.95 falls between the
exact ratio and its double. Wolfram evaluates, from each input's exact value at
60-digit precision, both the document's fair value and the pathwise value, the
latter as an expectation over a \texttt{GeometricDistribution} default month with the
payments' present values accumulated, and it proves symbolically that the
contractual payments discounted at the coupon are worth par.

% =========================================================================
\mkmvalidation
% =========================================================================

% Generated by tools.wolfram_oracle -- do not edit.
\begin{table}[H]\centering\small
\caption{Benchmark checks}\label{tab:bm-checks}
\begin{adjustbox}{max width=\textwidth}
\begin{tabular}{p{88mm}p{22mm}rl}
\toprule
Check & Result & Tolerance & Outcome \\
\midrule
Monthly payment (Eq.~1): model vs document & $2.0\times10^{-14}$ & $1.0\times10^{-12}$ & pass \\
Credit spread (Eqs.~6--13): model vs document & $6.4\times10^{-14}$ & $1.0\times10^{-12}$ & pass \\
PV of expected cashflows (Eqs.~19, 22): model vs document & $1.7\times10^{-14}$ & $1.0\times10^{-12}$ & pass \\
PV of expected losses (Eqs.~20, 23): model vs document & $3.8\times10^{-14}$ & $1.0\times10^{-12}$ & pass \\
Fair value (Eqs.~22, 24): model vs document & $1.7\times10^{-14}$ & $1.0\times10^{-12}$ & pass \\
Wolfram agrees with the reference (fair value and pathwise value) & $2.8\times10^{-30}$ & $1.0\times10^{-12}$ & pass \\
\bottomrule
\end{tabular}
\end{adjustbox}
\end{table}


\subsection{The correction}
% Generated by tools.wolfram_oracle -- do not edit.
\begin{table}[H]\centering\small
\caption{Fair value before (v1.1) and after (v1.3) the fix: GBP~300k property, 6\% coupon, 25 years, discounted at the coupon}\label{tab:bm-correction}
\begin{adjustbox}{max width=\textwidth}
\begin{tabular}{rrrrrrrr}
\toprule
LTV & $\delta$ & Income & $s$ (\%) & v1.1 & v1.3 & PV of loss & of par \\
\midrule
0.80 & 0.25 & 30{,}000 & 23.06 & 230{,}639 & 234{,}697 & 4{,}423 & +1.7\% \\
0.80 & 0.25 & 60{,}000 & 3.78 & 237{,}986 & 238{,}786 & 880 & +0.3\% \\
0.80 & 0.40 & 30{,}000 & 23.06 & 170{,}304 & 204{,}398 & 34{,}721 & +14.2\% \\
0.80 & 0.40 & 60{,}000 & 3.78 & 221{,}043 & 230{,}242 & 9{,}424 & +3.8\% \\
0.90 & 0.25 & 30{,}000 & 34.94 & 213{,}773 & 240{,}979 & 27{,}927 & +10.1\% \\
0.90 & 0.25 & 60{,}000 & 5.35 & 254{,}743 & 261{,}999 & 7{,}516 & +2.7\% \\
0.90 & 0.40 & 30{,}000 & 34.94 & 138{,}483 & 203{,}311 & 65{,}595 & +24.0\% \\
0.90 & 0.40 & 60{,}000 & 5.35 & 225{,}724 & 247{,}427 & 22{,}088 & +8.0\% \\
0.95 & 0.25 & 30{,}000 & 48.35 & 195{,}529 & 239{,}250 & 44{,}530 & +15.3\% \\
0.95 & 0.25 & 60{,}000 & 7.03 & 255{,}737 & 269{,}868 & 14{,}520 & +5.0\% \\
0.95 & 0.40 & 30{,}000 & 48.35 & 115{,}840 & 199{,}401 & 84{,}379 & +29.3\% \\
0.95 & 0.40 & 60{,}000 & 7.03 & 218{,}248 & 251{,}064 & 33{,}323 & +11.5\% \\
1.00 & 0.25 & 30{,}000 & 59.06 & 179{,}117 & 238{,}469 & 60{,}211 & +19.8\% \\
1.00 & 0.25 & 60{,}000 & 9.17 & 250{,}966 & 274{,}841 & 24{,}406 & +8.0\% \\
1.00 & 0.40 & 30{,}000 & 59.06 & 97{,}635 & 197{,}727 & 100{,}953 & +33.4\% \\
1.00 & 0.40 & 60{,}000 & 9.17 & 204{,}922 & 251{,}767 & 47{,}480 & +15.6\% \\
\bottomrule
\end{tabular}
\end{adjustbox}
\end{table}


\subsection{The model document's claims}
% Generated by tools.wolfram_oracle -- do not edit.
\begin{table}[H]\centering\small
\caption{The model document's claims, tested}\label{tab:bm-claims}
\begin{adjustbox}{max width=\textwidth}
\begin{tabular}{p{26mm}p{48mm}p{62mm}l}
\toprule
Claim & Stated & Benchmark finds & Outcome \\
\midrule
\S1, Eqs.~19, 24 & fair value is the PV of expected cashflows net of expected losses & the loss is deducted twice: fair value understated by up to 34.8\% of par & \textbf{not met} \\
Eq.~19 recovery & $B_{i-1}-\mathrm{LGD}_i$ & exceeds the recoverable value by a month's principal; up to 0.97\% of par & finding \\
\S6, Table~7 & spread and value by flood category & reproduced to the figures quoted; no loss given default at the baseline & met \\
\S6.1 & $+100$\,bp raises the 25-year payment 6--8\% & 8.3--11.9\% over rates of 2--9\% & \textbf{not met} \\
\S3.2.4 & zero income gives a 15\% spread & income is floored at GBP~1 first; the spread is 37.62\% & \textbf{not met} \\
\S7.2.2 & $\sum_i \mathrm{EL}_i\le B_0$ & at most 36.7\% of $B_0$ & met \\
\S5.3, \S7.2.1 & at the contractual rate the payments are worth par & identity proved symbolically & met \\
\S3.8 & discount at the contractual rate & production discounts at the risk-free rate; fair value moves up to 25\% & finding \\
\bottomrule
\end{tabular}
\end{adjustbox}
\end{table}


% =========================================================================
\mkmlimitations
% =========================================================================

\begin{enumerate}[leftmargin=*]
\item \textbf{Benchmark, not validation.} The code is shown to compute its documented
  formulae, and the fair value to be the expectation of the lender's cashflows under
  the model's default process; whether an affordability-based spread is the right
  default process is a separate question.
\item \textbf{Downstream.} Every fair value with a loss given default rises, and with
  it the portfolio totals and the discount to par that consume it; figures produced
  before the fix are understated by up to the amounts above.
\item \textbf{Discounting.} At the risk-free rate the loan calculator supplies, a
  coupon above it values the loan above par, so the discount to par measures the
  coupon's margin as well as credit (v1.3, \S5.3).
\end{enumerate}

\mkmhistorypage{27-Oct-2026 & 0.1 & First issue: tests v1.1; finds the loss counted twice & David K Kelly \\
27-Oct-2026 & 0.2 & After the fix: tests v1.3 & David K Kelly \\}

\end{document}
